\begin{table}[H]
    \centering
    \caption{Pemetaan Kebutuhan Fungsional ke Use Case}
    \label{tbl:mapping_fr_uc}
    \begin{tabular}{|p{0.10\textwidth}|p{0.35\textwidth}|p{0.45\textwidth}|}
        \hline
        \textbf{ID UC} & \textbf{Nama Use Case} & \textbf{Kebutuhan Fungsional (FR)} \\
        \hline
        UC-01 & Kelola Akun & \textbf{FR-01} (Autentikasi) \\
        \hline
        UC-02 & Mengerjakan Latihan Adaptif & \textbf{FR-02} (Sajian Materi), \textbf{FR-08} (Rekomendasi Konten) \\
        \hline
        UC-03 & Melakukan \textit{Peer Assessment} & \textbf{FR-03} (Memberi Feedback), \textbf{FR-09} (Intervensi Sosial) \\
        \hline
        UC-04 & Melihat Umpan Balik Rekan & \textbf{FR-03} (Menerima Feedback) \\
        \hline
        UC-05 & Memantau Profil Kompetensi & \textbf{FR-05} (Elo Personal), \textbf{FR-06} (Elo Sosial), \textbf{FR-07} (Model Terpadu) \\
        \hline
        UC-06 & Memantau Konteks Sosial & \textbf{FR-04} (Visualisasi Kelas) \\
        \hline
    \end{tabular}
\end{table}